\Qhead{2.3}{Свойства электромагнитных волн в анизотропных средах}{Свойства электромагнитных волн в анизотропных средах. Оптические кристаллы, уравнение Френеля, обыкновенная и необыкновенная волны. Магнитоактивные среды. Тензор диэлектрической проницаемости плазмы в магнитном поле; нормальные волны, их поляризация.}
\Subfile{2.3.1}{sections/02_waves/q03/s01}
\Subfile{2.3.2}{sections/02_waves/q03/s02}

\recap{В анизотропной среде ищут собственные поляризации тензорного волнового уравнения. В кристалле это обычная и необыкновенная волны, в продольно намагниченной плазме — круговые. Энергия и фаза могут распространяться в разных направлениях.}
\board{Матрица Френеля; сфера и эллипсоид; тензор $S,D,P$; два круговых показателя.}
\oral{Почему условие $\bm k\cdot\bm D=0$ не означает $\bm k\cdot\bm E=0$?}
\bait{Можно оставить снос необыкновенного луча: направление энергии задаёт нормаль к изочастотной поверхности, а не её радиус.}
\sources{Цысарь~\cite{tsysar1}, PDF-с.~56--60; нормальные волны: с.~58. Тензор холодной плазмы получен решением указанного уравнения движения; дополнительная проверка соглашений — Fitzpatrick~\cite{fitzpatrick}, раздел Cold-Plasma Dielectric Permittivity.}
